Language-model systems can answer from parametric weights, retrieve external evidence, read and write reusable memory, apply temporary task-time updates, or consolidate stable updates more durably. We study the controller problem: when should one policy choose among these six actions under a reward that combines answer quality, operation cost, forgetting risk, and delayed utility? We evaluate a unified router on frozen official runs covering stable recurrence, public changing-answer evaluation, edit propagation, and broad editing. On official PopQA recurring-memory, the router matches retrieval quality at 0.290 while reducing retrieval calls from 100.0 to 72.2; removing delayed utility drops quality to 0.090. On official FreshQA public main and sequence, retrieval is perfect across all four models, while suppressing delayed utility remains near ceiling at 0.982--0.996 and soft calibration only partially repairs the full router. On official MQuAKE, retrieval again dominates, and on official UniEdit all modes are near ceiling and largely nondiagnostic. These results support a simple empirical law: delayed utility helps when recurrence is stable and support can be reused safely, but retrieval-dominant behavior is more reliable in volatile public-answer and propagation-heavy regimes.
